% ==============================================================================
% chapters/ch04_plant_diagnostics.tex — บทที่ 4
% ==============================================================================
\chapter{การวินิจฉัยโรคพืช แมลงศัตรู และธาตุอาหาร}
\label{ch:plant_diagnostics}

การวินิจฉัยความผิดปกติของพืชได้อย่างรวดเร็วและแม่นยำเป็นหัวใจสำคัญของการทำเกษตรอัจฉริยะ ความเสียหายของพืชสามารถจำแนกออกเป็น 3 สาเหตุหลัก ได้แก่ โรคพืชที่เกิดจากเชื้อก่อโรค (ชีวปัจจัย: ชีววิทยาของเชื้อรา แบคทีเรีย ไวรัส), แมลงศัตรูพืชเข้าทำลาย และความผิดปกติทางสรีรวิทยาจากการขาดธาตุอาหาร (อชีวปัจจัย) ในบทนี้จะอธิบายรายละเอียดของชุดข้อมูลจริงที่ระบบรวบรวมไว้สำหรับพืชเศรษฐกิจสำคัญ

\academicsection{การวินิจฉัยโรคและแมลงศัตรูทุเรียน}
\label{sec:durian_diagnostics}

ทุเรียนเป็นราชาแห่งผลไม้และเป็นพืชเศรษฐกิจหลักของภาคตะวันออก อย่างไรก็ตาม ทุเรียนเป็นพืชที่อ่อนไหวต่อโรคเชื้อราและแมลงศัตรูพืชอย่างมาก ระบบได้คัดสรรภาพถ่ายจริงจำนวน 480 ภาพ ครอบคลุมความผิดปกติสำคัญดังนี้

\begin{enumerate}
    \item \textbf{โรคแอนแทรคโนสใบไหม้ (Leaf Blight / Colletotrichum):} เกิดจากเชื้อรา \textit{Colletotrichum} แผลมีลักษณะสีน้ำตาลเข้ม ไหม้ลามจากขอบใบหรือปลายใบเข้าสู่กลางใบ
    \item \textbf{โรคจุดสนิมใบจุดสาหร่าย (Algal Leaf Spot):} เกิดจากสาหร่ายสีเขียว \textit{Cephaleuros virescens} จุดแผลกลมนูน สีเขียวแกมส้มหรือสีสนิมเหล็ก มักระบาดในสวนทุเรียนที่มีความชื้นสูง
    \item \textbf{โรคใบจุดโฟมอปซิส (Phomopsis Leaf Spot):} แผลจุดสีน้ำตาลกระจายทั่วใบ ทำให้ใบร่วงก่อนกำหนด ส่งผลต่อการสังเคราะห์แสง
    \item \textbf{เพลี้ยไก่แจ้ทุเรียน (Durian Psyllid / Allocaridara attack):} ตัวอ่อนและตัวเต็มวัยของเพลี้ยไก่แจ้ \textit{Allocaridara malayensis} ดูดกินน้ำเลี้ยงที่ใบอ่อน ทำให้ใบหงิกงอ มีคราบมูลสีขาว และใบไหม้แห้งร่วงหล่น
\end{enumerate}

\academicsection{การวินิจฉัยโรคและแมลงศัตรูข้าว}
\label{sec:rice_diagnostics}

ข้าวเป็นพืชอาหารหลักของประเทศและประชากรโลก ในระบบได้นำเข้าภาพจริงจาก Rice Leaf Diseases Dataset และ Pest24 Dataset รวมกว่า 307 ภาพ ครอบคลุม:

\begin{itemize}
    \item \textbf{โรคขอบใบแห้ง (Bacterial Leaf Blight: BLB):} เกิดจากแบคทีเรีย \textit{Xanthomonas oryzae} แผลเริ่มจากขอบใบแห้งเป็นทางยาวสีเหลืองซีดและเปลี่ยนเป็นสีฟางข้าว
    \item \textbf{โรคใบจุดสีน้ำตาล (Brown Spot):} เกิดจากเชื้อรา \textit{Bipolaris oryzae} แผลรูปไข่สีน้ำตาลกระจายบนใบ มักสัมพันธ์กับดินที่ขาดธาตุอาหาร
    \item \textbf{โรคไหม้ข้าว (Rice Blast):} เกิดจากเชื้อรา \textit{Pyricularia oryzae} แผลรูปตาหรือรูปกระสวย ตรงกลางสีเทา ขอบสีน้ำตาลเข้ม
    \item \textbf{เพลี้ยกระโดดสีน้ำตาล (Brown Planthopper):} แมลงปากดูดที่สำคัญที่สุด ทำลายกอข้าวเกิดอาการแห้งตายทั้งแปลง (Hopperburn)
    \item \textbf{หนอนกอข้าว (Rice Stem Borer) และหนอนห่อใบ (Leaf Roller):} หนอนเจาะทำลายลำต้นและพับห่อใบข้าว
\end{itemize}

\academicsection{การตรวจวิเคราะห์โรคในข้าวโพดและมะเขือเทศ}
\label{sec:maize_tomato_diagnostics}

สำหรับข้าวโพดและมะเขือเทศ ระบบบรรจุภาพถ่ายจริงรวมกว่า 600 ภาพ จากคลังข้อมูลสากล PlantVillage, PlantDoc และ Corn Disease Benchmark ดังแสดงสรุปในตารางที่ \ref{tab:crops_disease_benchmark}

\begin{table}[htbp]
\caption{รายการโรคและแมลงศัตรูพืชในชุดข้อมูลจริงของระบบ}
\label{tab:crops_disease_benchmark}
\centering
\fontsize{9}{12}\selectfont
\begin{tabularx}{\textwidth}{llcX}
\toprule
\textbf{พืช} & \textbf{ชื่อโรค / แมลงศัตรู} & \textbf{จำนวนภาพจริง} & \textbf{ลักษณะอาการทางกายภาพ} \\
\midrule
ทุเรียน & แอนแทรคโนสใบไหม้ & 80 ภาพ & แผลไหม้สีน้ำตาลเข้มลามจากปลายใบ \\
ทุเรียน & จุดสนิมสาหร่าย & 80 ภาพ & จุดกลมนูนสีส้มคล้ายสนิมเหล็กบนหน้าใบ \\
ทุเรียน & เพลี้ยไก่แจ้ (Allocaridara) & 80 ภาพ & ใบหงิกงอ มีคราบมูลแป้งสีขาว \\
ทุเรียน & ใบสมบูรณ์ (Healthy) & 80 ภาพ & ใบสีเขียวเข้ม ผิวมันวาว ไร้ร่องรอยทำลาย \\
ข้าว & ขอบใบแห้ง (BLB) & 40 ภาพ & ขอบใบแห้งเป็นทางยาวขนานเส้นใบ \\
ข้าว & ใบจุดสีน้ำตาล (Brown Spot) & 40 ภาพ & จุดกลมรีสีน้ำตาล ขอบเขตชัดเจน \\
ข้าว & เพลี้ยกระโดดสีน้ำตาล & 60 ภาพ & ตัวเต็มวัยสีน้ำตาลดูดกินน้ำเลี้ยงกอข้าว \\
ข้าวโพด & ราสนิม (Common Rust) & 75 ภาพ & ตุ่มนูนสปอร์สีน้ำตาลแดงใต้ใบ \\
ข้าวโพด & ใบจุดสีเทา (Gray Leaf Spot) & 75 ภาพ & แผลรูปสี่เหลี่ยมผืนผ้าสีเทาระหว่างเส้นใบ \\
มะเขือเทศ & ใบจุดวงแหวน (Early Blight) & 75 ภาพ & แผลกลมสีน้ำตาลมีวงซ้อนศูนย์กลาง \\
มะเขือเทศ & ใบไหม้ (Late Blight) & 75 ภาพ & แผลไหม้สีเขียวช้ำอมน้ำ ลามเร็ว \\
\bottomrule
\end{tabularx}
\end{table}

ข้อมูลเหล่านี้ช่วยให้ผู้เข้าร่วมอบรมสามารถเปรียบเทียบอาการของโรคระหว่างพืชต่างชนิด และศึกษาผลกระทบของมุมมอง แสงเงา และสีฉากหลังต่อการทำงานของอัลกอริทึมปัญญาประดิษฐ์ได้อย่างครอบคลุม
